% GENERATED FILE. Edit canonical lessons, then run npm run generate:books.
% canonical-source-hash: fnv1a64:34a1de6bb420a131
% canonical-lessons: HI-C225-safkarna, HI-C225-jharna, HI-C225-pochna, HI-C225-nicorna, HI-C225-sukhana

\chapter{To clean, To dust, To wipe, To wring, To dry}
\label{ch:hi-to-clean-to-dust-to-wipe-to-wring-to-dry}
\hlchaptermodality{225}

\begin{chapteropening}
\textbf{By the end of this chapter:} \emph{I can say to clean, to dust, to wipe, to wring and to dry.}

Complete the last lesson of chapter 225: i can say to clean, to dust, to wipe, to wring and to dry.
\end{chapteropening}

\section[\texorpdfstring{\hi{साफ़} \hi{करना} (sāf karnā)}{साफ़ करना (sāf karnā)}]{\hi{साफ़} \hi{करना} (sāf karnā) — to clean}

\label{lesson:HI-C225-safkarna}



\begin{quote}
Before the new one: say the Hindi for a team, then the Hindi for a race, running.
\end{quote}

\subsection*{You'll want to know: \hi{साफ़} \hi{करना}}
\textbf{\hi{साफ़} \hi{करना}} — \emph{sāf karnā} — \textquotedblleft{}to clean\textquotedblright{}.

\begin{grammarlens}[title={What you've built}]
One more word for this chapter.
\end{grammarlens}

\subsection*{Your turn}
\begin{itemize}
\raggedright
  \item \emph{Say it:} \emph{sāf karnā}
  \item \emph{Say it:} \emph{sāf karnā}, once more
  \item \emph{Say it:} say \emph{sāf karnā}
  \item \emph{Recall:} say the Hindi for a team, then the Hindi for a race, running, then say \emph{sāf karnā} again
\end{itemize}

\begin{tcolorbox}[breakable,colback=teal!4,colframe=teal!35!black,title={Before you move on}]
What does \hi{साफ़} \hi{करना} mean? (\textquotedblleft{}To clean\textquotedblright{}.) Say it once more.
\end{tcolorbox}
\section[\texorpdfstring{\hi{झाड़ना} (jhāṛnā)}{झाड़ना (jhāṛnā)}]{\hi{झाड़ना} (jhāṛnā) — to dust, to sweep}

\label{lesson:HI-C225-jharna}



\begin{quote}
Before the new one: say the Hindi for a race, running, then the Hindi for to clean.
\end{quote}

\subsection*{You'll want to know: \hi{झाड़ना}}
\textbf{\hi{झाड़ना}} — \emph{jhāṛnā} — \textquotedblleft{}to dust, to sweep\textquotedblright{}.

\begin{grammarlens}[title={What you've built}]
One more word for this chapter.
\end{grammarlens}

\subsection*{Your turn}
\begin{itemize}
\raggedright
  \item \emph{Say it:} \emph{jhāṛnā}
  \item \emph{Say it:} \emph{jhāṛnā}, once more
  \item \emph{Say it:} say \emph{jhāṛnā}
  \item \emph{Recall:} say the Hindi for a race, running, then the Hindi for to clean, then say \emph{jhāṛnā} again
\end{itemize}

\begin{tcolorbox}[breakable,colback=teal!4,colframe=teal!35!black,title={Before you move on}]
What does \hi{झाड़ना} mean? (\textquotedblleft{}To dust, to sweep\textquotedblright{}.) Say it once more.
\end{tcolorbox}
\section[\texorpdfstring{\hi{पोंछना} (põchnā)}{पोंछना (põchnā)}]{\hi{पोंछना} (põchnā) — to wipe}

\label{lesson:HI-C225-pochna}



\begin{quote}
Before the new one: say the Hindi for to clean, then the Hindi for to dust.
\end{quote}

\subsection*{You'll want to know: \hi{पोंछना}}
\textbf{\hi{पोंछना}} — \emph{põchnā} — \textquotedblleft{}to wipe\textquotedblright{}.

\begin{grammarlens}[title={What you've built}]
One more word for this chapter.
\end{grammarlens}

\subsection*{Your turn}
\begin{itemize}
\raggedright
  \item \emph{Say it:} \emph{põchnā}
  \item \emph{Say it:} \emph{põchnā}, once more
  \item \emph{Say it:} say \emph{põchnā}
  \item \emph{Recall:} say the Hindi for to clean, then the Hindi for to dust, then say \emph{põchnā} again
\end{itemize}

\begin{tcolorbox}[breakable,colback=teal!4,colframe=teal!35!black,title={Before you move on}]
What does \hi{पोंछना} mean? (\textquotedblleft{}To wipe\textquotedblright{}.) Say it once more.
\end{tcolorbox}
\section[\texorpdfstring{\hi{निचोड़ना} (nicoṛnā)}{निचोड़ना (nicoṛnā)}]{\hi{निचोड़ना} (nicoṛnā) — to wring}

\label{lesson:HI-C225-nicorna}



\begin{quote}
Before the new one: say the Hindi for to dust, then the Hindi for to wipe.
\end{quote}

\subsection*{You'll want to know: \hi{निचोड़ना}}
\textbf{\hi{निचोड़ना}} — \emph{nicoṛnā} — \textquotedblleft{}to wring\textquotedblright{}.

\begin{grammarlens}[title={What you've built}]
One more word for this chapter.
\end{grammarlens}

\subsection*{Your turn}
\begin{itemize}
\raggedright
  \item \emph{Say it:} \emph{nicoṛnā}
  \item \emph{Say it:} \emph{nicoṛnā}, once more
  \item \emph{Say it:} say \emph{nicoṛnā}
  \item \emph{Recall:} say the Hindi for to dust, then the Hindi for to wipe, then say \emph{nicoṛnā} again
\end{itemize}

\begin{tcolorbox}[breakable,colback=teal!4,colframe=teal!35!black,title={Before you move on}]
What does \hi{निचोड़ना} mean? (\textquotedblleft{}To wring\textquotedblright{}.) Say it once more.
\end{tcolorbox}
\section[\texorpdfstring{\hi{सुखाना} (sukhānā)}{सुखाना (sukhānā)}]{\hi{सुखाना} (sukhānā) — to dry (something)}

\label{lesson:HI-C225-sukhana}



\begin{quote}
Before the new one: say the Hindi for to wipe, then the Hindi for to wring.
\end{quote}

\subsection*{You'll want to know: \hi{सुखाना}}
\textbf{\hi{सुखाना}} — \emph{sukhānā} — \textquotedblleft{}to dry (something)\textquotedblright{}.

\begin{grammarlens}[title={What you've built}]
One more word for this chapter.
\end{grammarlens}

\subsection*{Your turn}
\begin{itemize}
\raggedright
  \item \emph{Say it:} \emph{sukhānā}
  \item \emph{Say it:} \emph{sukhānā}, once more
  \item \emph{Say it:} say \emph{sukhānā}
  \item \emph{Recall:} say the Hindi for to wipe, then the Hindi for to wring, then say \emph{sukhānā} again
\end{itemize}

\begin{tcolorbox}[breakable,colback=teal!4,colframe=teal!35!black,title={Before you move on}]
What does \hi{सुखाना} mean? (\textquotedblleft{}To dry (something)\textquotedblright{}.) Say it once more.
\end{tcolorbox}
